\chapter*{Preface}
\addcontentsline{toc}{chapter}{Preface}

Book 11's Chapter 7 previewed an adjunction between formation and dissolution functors, establishing only enough structure to motivate the categorical vocabulary that book was building. Book 3's graded teleological operator $\mathcal{T}$ and Book 5's teleo-BV equation $Q_{\text{BV}}\mathcal{T} + \mathcal{T} Q_{\text{BV}} = 0$ both anticipated, in algebraic vocabulary, exactly the structure this book now develops categorically. The present volume elevates teleology, the purposive structure underlying this series' entropic smoothing since Book 3, into the full language of category theory: intention becomes a left adjoint, realization its right adjoint, and meaning the natural isomorphism connecting them.

This book's thesis is stricter than merely expressing purposive behavior in categorical vocabulary, and it is worth stating in advance, since every chapter that follows is an elaboration of it. Teleology is not the action of the future upon the present. It is the appearance, within forward dynamics, of constraints represented contravariantly and coupled to those dynamics by an adjunction. A terminal condition, an objective, an admissibility criterion, or a semantic target does not reach backward through time to shape the present that precedes it; it induces a contravariant map over the space of possibilities, and it is the adjunction between that contravariant map and the forward process it constrains, not any literal backward causal influence, that produces everything this series has previously described, in less precise vocabulary, as purpose. Where earlier volumes spoke of a system anticipating its own future, or of retrocausal structure compatible with monotonic entropy, this book replaces that vocabulary with its exact categorical content: structural coupling between a forward functor and a backward one, nothing more occult than that, and nothing less rigorous either. This is a categorical physics of purpose, closing the gap between thermodynamics, cognition, and ethics that every earlier volume in this series has approached from its own separate direction.

\chapter*{Diagrammatic Overview}
\addcontentsline{toc}{chapter}{Diagrammatic Overview}

The adjoint pair of functors $(\mathcal{A}, \mathcal{P})$ developed across this book, action and prediction, and the natural isomorphisms relating them, give teleology its categorical foundation directly: wherever this series has previously spoken informally of a system anticipating or acting purposively, the present book supplies the specific hom-set isomorphism that claim amounts to.

\chapter*{Reading Note}
\addcontentsline{toc}{chapter}{Reading Note}

Think of teleology not as goal-seeking but as categorical self-consistency. A system pursuing a goal, in the ordinary sense of that phrase, is on this book's account better described as a system whose action and prediction functors stand in a specific adjoint relationship, and the chapters below develop that relationship rather than the more familiar goal-seeking vocabulary the reader may be tempted to substitute for it.

\part{Foundations of Adjoint Teleodynamics}

\chapter{The Predictive and Active Functors}

Two endofunctors on $\mathbf{RSVPField}$ organize this book: $\mathcal{P}(X)$, the anticipated state a given object's dynamics point toward, and $\mathcal{A}(X)$, the realized state that dynamics actually produces. Both functors act on objects while preserving the thermodynamic structure Book 11 built into $\mathbf{RSVPField}$'s morphisms, so that prediction and action, whatever else distinguishes them, both remain curvature-preserving operations in the sense that book established.

The interpretation offered here treats prediction and action as dual curvature processes, with an example drawn from cognitive inference and behavioral enactment: a predicting system computes $\mathcal{P}(X)$ from its current state, while an acting system realizes $\mathcal{A}(X)$, and the entire teleological apparatus this book develops concerns the relationship between these two computations rather than treating either in isolation.

